\section{五大原则的综合框架}

将文章的五大原则整合来看，它们共同构成了一个完整的工具设计框架：

\begin{table}[H]
\centering
\caption{五大原则一览}
\begin{tabular}{clp{0.55\textwidth}}
\toprule
\textbf{编号} & \textbf{原则} & \textbf{核心要点} \\
\midrule
1 & 选择正确的工具集 & 少而精，面向工作流，整合多步操作 \\
2 & 工具命名空间 & 按服务和资源分层命名，帮助 Agent 快速定位 \\
3 & 返回有意义的上下文 & 高信号、语义化，避免低级标识符 \\
4 & 优化 Token 效率 & 分页、过滤、截断，引导高效行为 \\
5 & Prompt Engineering 描述 & 显式化隐含知识，消除歧义，结合评估迭代 \\
\bottomrule
\end{tabular}
\end{table}

\begin{warningbox}{各原则之间的相互作用}
五大原则并非独立运作，而是相互强化的。例如：选择了正确的工具集（原则 1）会减少命名空间的复杂度（原则 2）；优化返回值（原则 3）本身就是 Token 效率优化（原则 4）的一部分；而所有这些决策都需要通过 Prompt Engineering（原则 5）来正确传达给 Agent。如果只优化其中一个维度而忽略其他，整体效果可能会打折扣。
\end{warningbox}

这些原则的共同主题是：\textbf{以 Agent 的认知模型为中心进行设计}，而不是以底层系统的技术架构为中心。Agent 的上下文窗口是稀缺资源，每一个设计决策都应该围绕"如何让 Agent 用最少的上下文做最多的事"来展开。


